\section{Methodology: economics, mathematics, data science}
\label{sec:methodology}

This section makes the project's analytical contract precise: the event-driven
backtest semantics, the KPI formulas with their per-asset-class annualisation
factor, the strategy taxonomy, and the LLM-provider abstraction.

\subsection{Event-driven backtest semantics and look-ahead bias}
\label{sec:lookahead}

Look-ahead bias is the canonical pitfall of backtest construction: a strategy
``sees'' information that, in real time, it would not have had access to when it
needed to act \cite{lopezdeprado2018,bailey2014}. The most common form is a
\emph{single-bar} leak \textemdash{} a strategy reads bar~$t$'s closing price and
then places an order that fills at bar~$t$'s open or close. Across a long backtest
the false edge accumulates and inflates every downstream metric. QuantEdge
forbids the leak architecturally: orders generated during bar~$t$ fill at the
\emph{open of bar~$t{+}1$}, with commission and slippage applied to the fill
price. The rule is implemented once, in \texttt{backend/backtest/engine.py}, and
asserted by an oracle-strategy test in \texttt{tests/test\_engine\_no\_lookahead.py}
that injects a strategy that ``knows'' the future close; the test passes only if
the engine still forces the oracle to trade at next-bar open.

In pseudocode the engine loop is:

\begin{lstlisting}[language=Python]
for t in range(len(bars)):
    # Target position from data up to and including bar t.
    target = strategy.generate_signals(bars[:t+1]).iloc[t]   # in {-1, 0, 1}
    if t == len(bars) - 1: break       # no t+1 open to fill at
    delta = target - portfolio.current_position
    if delta == 0: continue
    # Fill at bar (t+1) open, applying slippage and commission.
    fill_price = bars.iloc[t+1].open * (1 + sign(delta) * slippage_bps / 10_000)
    commission = abs(delta * portfolio.size) * fill_price * commission_bps / 10_000
    portfolio.apply(delta, fill_price, commission)
\end{lstlisting}

The signal contract is uniformly \texttt{pd.Series[int]} with values in
$\{-1, 0, +1\}$ for ``short, flat, long''. Strategies do not multiply their own
signals by a position-size factor; the engine handles sizing via either
\emph{fixed fractional} (a fraction of equity per trade) or \emph{volatility
targeting} (position scaled inversely to realised volatility). An optional
circuit-breaker forces the run flat if equity draws down beyond a user-supplied
threshold.

\subsection{Performance metrics}
\label{sec:kpi}

The KPI definitions follow the conventions in Bacon~\cite{bacon2008}. Let
$\{r_t\}_{t=1}^{T}$ denote the per-bar returns of the equity curve
$\{V_t\}_{t=0}^{T}$, $k$ the number of bars in one calendar year
(Section~\ref{sec:calendars}), and $y = T/k$ the run's length in years.

\paragraph{Total return and CAGR.}
\begin{equation}
  R = \frac{V_T}{V_0} - 1,\qquad
  \mathrm{CAGR} = \left(\frac{V_T}{V_0}\right)^{1/y} - 1.
\end{equation}

\paragraph{Annualised volatility.}
\begin{equation}
  \sigma_{\text{ann}} = \sqrt{k}\,\sigma(r_t).
\end{equation}

\paragraph{Sharpe and Sortino ratios.}
\begin{equation}
  \mathrm{Sharpe} = \sqrt{k}\,\frac{\bar r}{\sigma(r_t)},\qquad
  \mathrm{Sortino} = \sqrt{k}\,\frac{\bar r}{\sigma_{-}(r_t)},
\end{equation}
where $\sigma_{-}(r_t)$ is the downside semi-deviation $\sqrt{\mathbb{E}[\min(r_t, 0)^2]}$.
We adopt the zero-risk-free-rate convention for retail backtesting
\cite{sharpe1966,sortino1994}.

\paragraph{Maximum drawdown and Calmar ratio.}
\begin{equation}
  \mathrm{maxDD} = \min_{t} \frac{V_t - \max_{s \le t} V_s}{\max_{s \le t} V_s},\qquad
  \mathrm{Calmar} = \frac{\mathrm{CAGR}}{|\mathrm{maxDD}|}.
\end{equation}
Calmar~\cite{young1991} captures the pain-adjusted return: how much annual growth
the strategy delivered per unit of drawdown the user had to endure.

\paragraph{Trade-level metrics.} Let $\{p_i\}_{i=1}^{N}$ denote per-leg net P\&L:
\begin{equation}
  \mathrm{winrate} = \frac{|\{i:p_i > 0\}|}{N},\qquad
  \mathrm{PF} = \frac{\sum_{p_i > 0} p_i}{\sum_{p_i < 0} |p_i|},
\end{equation}
\begin{equation}
  \mathrm{expectancy} = \mathrm{winrate}\cdot \overline{p_{+}} - (1-\mathrm{winrate}) \cdot |\overline{p_{-}}|,
\end{equation}
where $\overline{p_{+}}$ and $\overline{p_{-}}$ are the means of the positive and
negative P\&L populations respectively~\cite{pardo2008,tharp1998}.

\subsection{Annualisation factor by asset class and timeframe}
\label{sec:calendars}

A Sharpe ratio computed against the wrong $k$ misstates the annualised number by
tens of percent. QuantEdge looks up $k$ from
\texttt{backend/analytics/periods.py} as a function of \texttt{(asset\_class,
timeframe)}:

\begin{table}[h]
\centering
\small
\begin{tabular}{lcc}
\toprule
\textbf{Asset class} & \textbf{Daily $k$} & \textbf{Hourly $k$} \\
\midrule
Equity / ETF (NYSE) & 252  & 1\,638 \\
FX (24$\times$5)    & 260  & 6\,240 \\
Crypto (24$\times$7)& 365  & 8\,760 \\
\bottomrule
\end{tabular}
\caption{Annualisation factor $k$ per asset class and timeframe. NYSE hourly uses
$252 \times 6.5$ elapsed market hours per session.}
\end{table}

The NYSE hourly value reflects elapsed market hours, not the literal bar count,
because the first hourly bar of a regular session covers only the opening
30-minute auction. Early-close sessions (about ten per year) further reduce the
count slightly; the lookup is calibrated against the
\texttt{exchange\_calendars} XNYS schedule rather than computed analytically.

\subsection{Strategy taxonomy}

The eleven shipped strategies populate five families. Each family responds to a
different market regime; running them side-by-side makes that regime sensitivity
visible, which is the explicit pedagogical motive for shipping the breadth.

\begin{itemize}
  \item \textbf{Trend} (SMA Crossover \cite{pardo2008}, MACD Crossover \cite{appel2005},
        Ichimoku Cloud \cite{hosoda1969}): designed to capture sustained directional moves.
  \item \textbf{Momentum} (Time Series Momentum \cite{moskowitz2012}): own-asset trailing-return
        sign; profits from 1\textendash{}12 month persistence.
  \item \textbf{Breakout} (Donchian \cite{faith2007}, Keltner \cite{keltner1960}):
        long on channel breaks; trades clean trending regimes.
  \item \textbf{Mean reversion} (Bollinger \cite{bollinger2001}, CCI \cite{lambert1980},
        RSI \cite{wilder1978}, Stochastic Oscillator \cite{lane1984}): trade fades within
        range-bound regimes; bleed in trends.
  \item \textbf{Benchmark} (Buy and Hold): always long; serves as the reference both
        for active strategies and for the on-chart same-asset overlay.
\end{itemize}

All eleven strategies share the same signal contract and therefore the same engine
path. Adding a twelfth is a single-file change plus a registry edit; the frontend's
parameter form regenerates from the new JSON Schema.

\subsection{The LLM-provider abstraction}

Two of the six runtime agents \textemdash{} the Orchestrator and the Explanation
Agent \textemdash{} are LLM-backed, and every call they make goes through
\texttt{LLMProvider.generate(\ldots)}. Two implementations ship:
\texttt{NullProvider} (deterministic canned text) and \texttt{GeminiProvider}
(Google's Gemini API). The abstraction earns its keep three ways: tests use
\texttt{NullProvider} exclusively and stay reproducible without an API key;
switching provider is a config change, not a code change; and it makes the
project's policy explicit \textemdash{} \emph{the deterministic surface is the
default}, and any non-determinism is a conscious operator opt-in
(\texttt{LLM\_ENABLED=true} plus a populated \texttt{GEMINI\_API\_KEY}).
